% GENERATED by scripts/build-lecture-slides.py; do not edit. Portable profile: XeTeX, theme and fonts from the core.
\makeatletter\IfFileExists{core/beamer/beamerthememx.sty}{\def\mxcorepath{core/}\def\input@path{{core/beamer/}{core/}}}{\def\mxcorepath{../core/}\def\input@path{{../core/beamer/}{../core/}}}\makeatother
\documentclass[aspectratio=169,10pt,t]{beamer}
\usetheme{mx}
\usepackage{graphicx,booktabs}
\usepackage[french]{babel}
\setlength{\parskip}{.45em}
\hypersetup{pdfauthor={},pdfsubject={Lecture slides},pdfkeywords={},pdftitle={Trigonométrie}}
\providecommand{\N}{\mathbb{N}}
\providecommand{\Z}{\mathbb{Z}}
\providecommand{\D}{\mathbb{D}}
\providecommand{\Q}{\mathbb{Q}}
\providecommand{\R}{\mathbb{R}}
\providecommand{\C}{\mathbb{C}}
\providecommand{\vect}[1]{\overrightarrow{#1}}
\providecommand{\vi}{\vec{\imath}}
\providecommand{\vj}{\vec{\jmath}}
\providecommand{\vk}{\vec{k}}
\providecommand{\norm}[1]{\left\lVert #1\right\rVert}
\providecommand{\vc}[2]{\begin{pmatrix}#1\\#2\end{pmatrix}}
\providecommand{\deter}[2]{\det\left(#1,#2\right)}
\providecommand{\abs}[1]{\left\lvert #1\right\rvert}
\providecommand{\dd}{\,\mathrm{d}}
\providecommand{\e}{\mathrm{e}}
\providecommand{\Card}{\operatorname{Card}}
\providecommand{\E}{\operatorname{E}}
\providecommand{\V}{\operatorname{V}}
\providecommand{\degre}{^{\circ}}
\providecommand{\un}[1]{\,\mathrm{#1}}
\providecommand{\sci}[2]{#1\times10^{#2}}
\providecommand{\pgcd}{\operatorname{PGCD}}
\providecommand{\ppcm}{\operatorname{PPCM}}
\title{Trigonométrie}
\subtitle{Cours et exemples}
\date{}
\author{}
\begin{document}
\mxtitleframe
\begin{frame}{}
Sur des cercles de rayons différents, construire un arc de longueur égale au
rayon. Comparer l'angle au centre obtenu. Cette invariance conduit à définir le
radian; en déduire la mesure d'un tour complet.
\end{frame}
\begin{frame}{Objectifs}
\begin{itemize}

  \item Convertir entre degrés et radians.
  
\item Utiliser le cercle trigonométrique.
  
\item Connaître les valeurs remarquables de sin, cos, tan.
  
\item Utiliser \(\cos^2\theta + \sin^2\theta = 1\).
  
\item Résoudre graphiquement \(\cos x = a\) et \(\sin x = b\).

\end{itemize}
\end{frame}
\begin{frame}{Radians}
Le \textbf{radian} est l'unité d'angle SI. Un angle de \(1\) rad correspond
à un arc de longueur égale au rayon.
\[\alpha_\circ \times \dfrac{\pi}{180} = \alpha_\text{rad},\qquad
  \alpha_\text{rad} \times \dfrac{180}{\pi} = \alpha_\circ.\]
\end{frame}
\begin{frame}{Valeurs remarquables et relations}
\begin{center}\begin{tabular}{lllllll}
\(\theta\) & \(0\) & \(\dfrac{\pi}{6}\) & \(\dfrac{\pi}{4}\) & \(\dfrac{\pi}{3}\) & \(\dfrac{\pi}{2}\) & \(\pi\) \\ 
\(\cos\theta\) & \(1\) & \(\dfrac{\sqrt{3}}{2}\) & \(\dfrac{\sqrt{2}}{2}\) & \(\dfrac{1}{2}\) & \(0\) & \(-1\) \\ 
\(\sin\theta\) & \(0\) & \(\dfrac{1}{2}\) & \(\dfrac{\sqrt{2}}{2}\) & \(\dfrac{\sqrt{3}}{2}\) & \(1\) & \(0\) \\ 
\end{tabular}\end{center}
\end{frame}
\begin{frame}{Valeurs remarquables et relations}
Pour tout angle \(\theta\) :
\[\cos^2\theta + \sin^2\theta = 1.\]
\end{frame}
\begin{frame}{Le cercle trigonométrique}
On place un angle \(\theta\) (en radians) sur le cercle unité, en partant
du point \(A(1,0)\) et en tournant dans le sens direct (sens anti-horaire).
Le point \(M\) ainsi obtenu a pour coordonnées :
\[M(\cos\theta,\, \sin\theta).\]
Pour \(\cos\theta\ne0\), on définit aussi la \textbf{tangente} :
\[\tan\theta = \dfrac{\sin\theta}{\cos\theta}.\]
\end{frame}
\begin{frame}{Le cercle trigonométrique}
\begin{center}\begin{tabular}{lll}
Symétrie & Cosinus & Sinus \\ 
Opposé : \(-\theta\) & \(\cos(-\theta)=\cos\theta\) & \(\sin(-\theta)=-\sin\theta\) \\ 
Supplément : \(\pi-\theta\) & \(\cos(\pi-\theta)=-\cos\theta\) & \(\sin(\pi-\theta)=\sin\theta\) \\ 
Antipose : \(\pi+\theta\) & \(\cos(\pi+\theta)=-\cos\theta\) & \(\sin(\pi+\theta)=-\sin\theta\) \\ 
Période : \(\theta+2k\pi\) & \(\cos(\theta+2k\pi)=\cos\theta\) & \(\sin(\theta+2k\pi)=\sin\theta\) \\ 
\end{tabular}\end{center}
\end{frame}
\begin{frame}{Le cercle trigonométrique}
\begin{enumerate}
\setcounter{enumi}{0}
  \item Repérer le quadrant (\(\cos\) et \(\sin\) sont positifs ou négatifs).
  
\item Identifier l'angle de référence (angle avec l'axe des \(x\)).
  
\item Appliquer la symétrie correspondante.

\end{enumerate}

\emph{Exemple :} \(\cos(\dfrac{5\pi}{6})\).
\(\dfrac{5\pi}{6} = \pi - \dfrac{\pi}{6}\) : second quadrant.
\(\cos(\dfrac{5\pi}{6}) = -\cos(\dfrac{\pi}{6}) = -\dfrac{\sqrt{3}}{2}\).
\end{frame}
\begin{frame}{Résolution d'équations trigonométriques}
\begin{enumerate}
\setcounter{enumi}{0}
  \item Si \(|a|>1\) : pas de solution.
  
\item Trouver l'angle de référence \(\alpha\in[0,\pi]\) tel que \(\cos\alpha=a\).
  
\item Les solutions sur \([0,2\pi]\) sont : \(x=\alpha\) et \(x=2\pi-\alpha\).
        (Symétrie par rapport à l'axe \(Ox\).)

\end{enumerate}
\end{frame}
\begin{frame}{Résolution d'équations trigonométriques}
\begin{enumerate}
\setcounter{enumi}{0}
  \item Si \(|a|>1\) : pas de solution.
  
\item Trouver l'angle de référence \(\alpha\in[-\dfrac{\pi}{2},\dfrac{\pi}{2}]\) tel que \(\sin\alpha=a\).
  
\item Les solutions sur \([0,2\pi]\) sont : \(x=\alpha\) (adapté à \([0,2\pi]\))
        et \(x=\pi-\alpha\) (adapté à \([0,2\pi]\)).
        (Symétrie par rapport à l'axe \(Oy\).)

\end{enumerate}
\end{frame}
\begin{frame}{Résolution d'équations trigonométriques}
L'angle de référence : \(\cos(\dfrac{\pi}{3})=\dfrac{1}{2}\), donc \(\cos(\pi-\dfrac{\pi}{3})=-\dfrac{1}{2}\).
\(\alpha = \dfrac{2\pi}{3}\).
Solutions : \(x = \dfrac{2\pi}{3}\) et \(x = 2\pi - \dfrac{2\pi}{3} = \dfrac{4\pi}{3}\).
\end{frame}
\begin{frame}{Résolution d'équations trigonométriques}
\(\sin(\dfrac{\pi}{4})=\dfrac{\sqrt{2}}{2}\).
Solutions : \(x = \dfrac{\pi}{4}\) et \(x = \pi - \dfrac{\pi}{4} = \dfrac{3\pi}{4}\).
\end{frame}
\begin{frame}{Résolution d'équations trigonométriques}
\(\tan\theta = \dfrac{\sin\theta}{\cos\theta}\) est définie pour \(\cos\theta\ne0\),
c'est-à-dire \(\theta\ne\dfrac{\pi}{2}+k\pi\).

Valeurs remarquables :
\(\tan(0)=0\), 
\(\tan(\dfrac{\pi}{6})=\dfrac{1}{\sqrt{3}}=\dfrac{\sqrt{3}}{3}\), 
\(\tan(\dfrac{\pi}{4})=1\), 
\(\tan(\dfrac{\pi}{3})=\sqrt{3}\).
\end{frame}
\end{document}
